\begin{center}\begin{tikzpicture}[font=\small]
\draw[black!30,rounded corners=2pt] (0,0) rectangle (14.4,8.2);
\node[anchor=west,font=\small\bfseries] at (.5,7.72) {消元不是另造点：两个实根分别对应原来的两个交点};
% MODEL 8 4 1 1 -1
\begin{scope}[shift={(7.2,3.65)},scale=1.08]
\coordinate (P) at (-0.8968052533,1.8968052533);
\coordinate (Q) at (2.2301385866,-1.2301385866);
\coordinate (M) at (0.6666666667,0.3333333333);
\coordinate (O) at (0.0000000000,0.0000000000);
\draw[blue,very thick] (0,0) ellipse (2.8284271247 and 2.0000000000);
\draw[-Stealth,black!45] (-3.3,0)--(3.3,0) node[right] {$x$};\draw[-Stealth,black!45] (0,-2.3)--(0,2.3) node[above] {$y$};
\draw[orange!85!black,very thick] (P)--(Q);\draw[green!45!black,dashed] (P)--(P|-O);\draw[green!45!black,dashed] (Q)--(Q|-O);\fill[black] (P) circle (1.8pt) node[above left=5pt] {$P(x_1,y_1)$};
\fill[black] (Q) circle (1.8pt) node[below right=5pt] {$Q(x_2,y_2)$};
\node[below=7pt,text=green!45!black] at (P|-O) {$x_1$};\node[above=7pt,text=green!45!black] at (Q|-O) {$x_2$};
\end{scope}
\node[anchor=west] at (.8,6.7) {$x^2/8+y^2/4=1,\quad x+y-1=0$};\node[anchor=west] at (.8,.6) {两交点的横坐标，正是消元方程的两个实根};
\end{tikzpicture}\end{center}
